\chapter*{Preface}
\addcontentsline{toc}{chapter}{Preface}

This book supports several styles of instruction. Followed sequentially, the
chapters form a conventional undergraduate trigonometry course, self-contained
even if the comparative sidebars are skipped for time. Read alongside the
derivations, they support a flipped classroom in which class time is reserved for
problem solving and proof construction. Read slowly, with the motivating
questions at the head of each section taken seriously before the formal
definitions that follow them, they support inquiry-based instruction.

Each section follows a consistent rhythm: a motivating problem, a derivation, a
worked example, a short practice-now check, and a set of exercises. Major
constructions are additionally given as explicit, numbered procedures, since
students often need a procedure to lean on before they can appreciate the theory
that generated it. Exercises are meant to be worked in order; later exercises
frequently depend on results established in earlier ones. Answers to selected
exercises appear in the appendix.

\bigskip
\noindent Flyxion \\
Independent Researcher
